\chapter{Cadrage du projet et cahier des charges}
\section{Finalité et périmètre}
Le projet s'inscrit dans la SAÉ « Autour du tracé d'un métro circulaire » de
Sup Galilée\cite{enonce}. Le sujet demande notamment de justifier la NP-difficulté,
de visualiser les données et les solutions, de programmer une heuristique et une
métaheuristique, d'implémenter une formulation compacte et de comparer les méthodes.
La version étudiée fournit en plus une recherche locale distincte de la recherche
tabou, des campagnes à plusieurs graines et des résultats consultables par des
résumés Markdown.

Un point représente une zone candidate à l'implantation d'une station. Les points
ont tous le même poids dans le coût d'accès : aucune population, demande ou
fréquence de déplacement n'est fournie. Les distances sont calculées à vol
d'oiseau dans le plan des coordonnées. Elles ne décrivent ni un chemin dans la
voirie ni une durée de transport. Le mot « coût » désigne ici la fonction objectif,
sans conversion en euros.

\section{Entrées et décisions attendues}
Les entrées sont un fichier de coordonnées, un nombre de stations $p$, un poids
$\alpha$ et les paramètres de la méthode. Le premier sommet du fichier est une
station imposée. Les bornes retenues sont $3\leq p\leq n$ et $0\leq\alpha\leq10$.
Le code utilise des indices internes de 0 à $n-1$ ; les exports lisibles numérotent
les points de 1 à $n$. Ainsi, le sommet imposé habituel est 0 en interne et 1
dans les résultats.

Une sortie complète contient l'ensemble des stations, l'ordre du cycle, une
affectation par point et les trois coûts : anneau non pondéré, accès et total
pondéré. Elle précise aussi la méthode, la graine, les budgets et le statut.
Pour les recherches, les traces distinguent le coût courant de la meilleure
valeur rencontrée. Pour CBC, la preuve d'optimalité et la borne inférieure sont
des informations séparées de la solution réalisable.

\section{Exigences de qualité}
\exigence{Réalisabilité}{Toute solution doit comporter exactement $p$ stations,
inclure le sommet imposé et former un cycle simple passant par chacune d'elles.
Chaque point est affecté à une station ouverte et chaque station à elle-même.}
\exigence{Cohérence numérique}{Les méthodes partagent la même matrice de distances
et la même évaluation du coût. Les coordonnées non finies et les paramètres
hors domaine sont rejetés avant l'optimisation.}
\exigence{Comparabilité}{Les méthodes approchées partent de la même construction
pour un cas fixé. Les comparaisons conservent l'instance, $p$ et $\alpha$ ; les
statistiques ne mélangent pas des problèmes différents.}
\exigence{Traçabilité}{Une figure est reliée à une solution sauvegardée. Les
coordonnées et l'ordre du cycle sont conservés pour régénérer la figure sans
recalculer un autre parcours.}
\exigence{Utilisation}{L'installation, les tests et les exécutions sont accessibles
par une interface de commande et un Makefile. Les résultats individuels et les
campagnes possèdent chacun un point d'entrée documenté.}

\section{État du travail étudié}
Le dépôt public est \href{https://github.com/birane-m/ring-star-problem}{birane-m/ring-star-problem}.
Les trois commits publiés jusqu'à \code{8761003} regroupent le code expérimental,
les outils du projet et la documentation. Les calculs présentés ont été enregistrés
le 17 octobre 2026 ; leur empreinte de sources est conservée dans les métadonnées.
Ils précèdent les dernières modifications de présentation et ne sont pas
présentés comme une campagne relancée après chaque commit.

Le contrôle documentaire a identifié une écriture finale manquante du CSV pour
un cas PLNE ignoré. La correction jointe aux sources ne modifie aucun algorithme
ni aucun coût ; elle fiabilise la complétude de l'export. Son statut et les
contrôles associés sont explicités au chapitre~\ref{chap:qualite}.
\bilan{Le projet répond à une question de localisation-routage dans un modèle
géométrique défini. La validité des comparaisons dépend autant des conventions
communes et des budgets que des valeurs finales des coûts.}
